\begin{resume}
\ifisanon
    \ifisblind
        \isanonfalse
    \else
        \ifismaster
        该论文作者在学期间取得的阶段性成果（学术论文等）已满足我校硕士学位评阅相关要求。为避免阶段性成果信息对专家评价学位论文本身造成干扰，特将论文作者的阶段性成果信息隐去。
        \else
        % 该论文作者在学期间取得的阶段性成果（学术论文等）已满足我校博士学位评阅相关要求。为避免阶段性成果信息对专家评价学位论文本身造成干扰，特将论文作者的阶段性成果信息隐去。
          \section*{发表的学术论文} % 发表的和录用的合在一起

          \begin{enumerate}[label={[\arabic*]},labelsep=6pt,topsep=0pt,partopsep=0pt,
        parsep=0pt,itemsep=1pt,leftmargin=0em,itemindent=42pt]
          % \item 第一作者.期刊论文.IEEE Transactions on Neural Networks and Learning Systems,2025,SCI影响因子8.9.
          % \item 第一作者.期刊论文.IEEE Transactions on Neural Systems and Rehabilitation Engineering,2024,SCI影响因子4.8.
          % \item 第一作者.期刊论文.IEEE Transactions on Cognitive and Developmental Systems,2023,SCI影响因子5.0.
          % \item 第一作者.会议论文.International Conference on Automation Control,Algorithm,and Intelligent Bionics,2024.
          % \item 第二作者.期刊论文.International Journal of Computer Vision,2024,SCI影响因子11.6.
          % \item 第二作者.期刊论文.Physics of Fluids,2024,SCI影响因子4.6.
          % \item 第四作者.期刊论文.IEEE Transactions on Medical Imaging,2025,SCI影响因子10.6.
          % \item 第五作者.期刊论文.Frontiers in Neurorobotics,2025,SCI影响因子2.9.
          % \item 第五作者.期刊论文.IEEE Transactions on Human-Machine Systems,2023,SCI影响因子3.5.
          % \item 第二作者.会议论文.Youth Academic Annual Conference of Chinese Association of Automation,2022.
          % \item 第三作者.会议论文.IEEE International Conference on Civil Aviation Safety and Information Technology,2022.
          % \item 第四作者.会议论文.Asian Conference on Artificial Intelligence Technology,2023.
          % \item 第二作者.会议论文.Asian Conference on Artificial Intelligence Technology,2025.
          \item 第一作者.DMAE-EEG: A Pretraining Framework for EEG Spatiotemporal Representation Learning.IEEE Transactions on Neural Networks and Learning Systems,2025. (SCI 收录, 中科院1区, DOI:10.1109/TNNLS.2025.3581991).
          \item 第一作者.MASER: Enhancing EEG Spatial Resolution With State Space Modeling.IEEE Transactions on Neural Systems and Rehabilitation Engineering,2024,32:3858-3868.(SCI 收录, 中科院1区, DOI:10.1109/TNSRE.2024.3481886).
          \item 第一作者.Graph learning with co-teaching for EEG-based motor imagery recognition.IEEE Transactions on Cognitive and Developmental Systems,2023,15:1722-1731. (SCI 收录, 中科院3区, DOI:10.1109/TCDS.2022.3174660).
          \item 第一作者.MAE-MI: masked autoencoder for self-supervised learning on motor imagery EEG data.International Conference on Automation Control, Algorithm, and Intelligent Bionics,2024,13259:805-810.(EI 收录, DOI:10.1117/12.3039558).
          \item 第二作者.SplatFlow: Learning multiframe optical flow via splatting.International Journal of Computer Vision,2024,132(8):3023-3045. (SCI 收录, 中科院2区, DOI:10.1007/ s11263-024-01993-0).
          \item 第二作者.Amphibious vehicle's resistance optimization through neural networks and genetic algorithms.Physics of Fluids,2024,36(6):065129. (SCI 收录, 中科院2区,DOI: 10.1063/5.0210244).
          \item 第四作者.M₂DC: A Meta Learning Framework for Generalizable Diagnostic Classification of Major Depressive Disorder.IEEE Transactions on Medical Imaging,2025, 44(2):855-867. (SCI 收录, 中科院1区,DOI:10.1109/TMI.2024.3461312).
          \item 第五作者.Universal slip detection of robotic hand with tactile sensing.Frontiers in Neurorobotics,2025,19:1478758. (SCI 收录, 中科院4区, DOI:10.3389/fnbot.2025.147 8758).
          \item 第五作者.Fusion of spatial, temporal, and spectral EEG signatures improves multilevel cognitive load prediction.IEEE ransactions on Human-Machine Systems,2023, 53(2):357-366. (SCI 收录, 中科院2区,DOI:10.1109/THMS.2023.3235003).
          \item 第二作者.Combining multi-scale convolutional neural network and transformers for EEG-based RSVP detection.2022 37th Youth Academic Annual Conference of Chinese Association of Automation,2022:426-431. (EI 收录,DOI:10.1109/YAC57282.2022. 10023688).
          \item 第三作者.Micro-expression Recognition Using Pretrained Model and Transformer.2022 IEEE 4th International Conference on Civil Aviation Safety and Information Technology,2022:1404-1408. (EI 收录, DOI:10.1109/ICCASIT55263.2022.9987236).
          \item 第四作者.Co-Teaching Learning from Noisy Labeled FMRI Data for Diagnostic Classification of Major Depression.2023 7th Asian Conference on Artificial Intelligence Technology,2023:404-409. (EI 收录, DOI:10.1109/ACAIT60137.2023.10528543).
          \item 第二作者.Efficient Transformer-based Motor Imagery Classification for Edge-deployed Brain-Computer Interfaces.Asian Conference on Artificial Intelligence Technology,2025. (EI 已录用).
          \item 第一作者.Co-Guidance Learning for EEG Channel Selection with State Space Modeling. (SCI 在审).
          \end{enumerate}

          \section*{申请的发明专利} % 有就写，没有就删除
          \begin{enumerate}[label={[\arabic*]},labelsep=6pt,topsep=0pt,partopsep=0pt,
        parsep=0pt,itemsep=1pt,leftmargin=0em,itemindent=42pt]
          \item 第二作者.一种基于图神经网络的多类别运动想象任务识别方法:中国,专利号：ZL202210228171.4.(已授权).
          \item 第二作者.一种数据驱动的脑电图通道扩展方法及系统:中国,专利号：ZL202411955772.0. (已授权).
          \item 第三作者.一种基于多脑BCI的目标增强判读方法和装置:中国,专利号：ZL202210228171.4.(已授权).
          \item 第二作者.一种面向多源脑电图信号的自监督预训练方法:中国,申请号：2023117321526.(实质性审查).
          \item 第二作者.一种基于注意力机制的缺损脑电图数据修复方法:中国,申请号：2024102739030.(实质性审查).
          \item  第二作者.一种基于脑电分析的认知状态评估方法与可视化反馈系统:中国,申请号：2025110578957.(实质性审查).
          \item 第三作者.一种无人值守系统及其基于光流的监控方法:中国,申请号：2022112118920.(实质性审查).
          \item 第二作者.一种基于双分支协同引导学习的EEG通道选择方法与系统:中国,申请号：2025110578590.(已受理).
          \end{enumerate}

          \section*{获得的竞赛奖项} %

        \begin{enumerate}[label={[\arabic*]},labelsep=6pt,topsep=0pt,partopsep=0pt,
        parsep=0pt,itemsep=1pt,leftmargin=0em,itemindent=42pt]
        \normalsize
        \item 2024年第二十八届全国发明展览会“一带一路”暨金砖国家技能发展与技术创新大赛,金奖.
        \item 2024年第六届“华为杯”中国研究生人工智能创新大赛,国家三等奖.
        \item 2023年全球人工智能技术创新大赛硬件挑战赛无人对抗组,国家二等奖.
        \item 2023年全球人工智能技术创新大赛硬件挑战赛深度学习组,国家一等奖.
        \item 2023年第九届“新动航空杯”中国研究生未来飞行器创新大赛,国家二等奖.
        \item 2023年第五届“申昊杯”中国研究生机器人创新设计大赛,国家二等奖.
        \item 2022年第四届“申昊杯”中国研究生机器人创新设计大赛,国家二等奖.
        \end{enumerate}

        \section*{参与的科研项目} %

          \begin{enumerate}[label={[\arabic*]},labelsep=6pt,topsep=0pt,partopsep=0pt,
          parsep=0pt,itemsep=1pt,leftmargin=0em,itemindent=42pt]
          \normalsize
          \item 侵入式/非侵入式双向闭环脑机接口关键技术及系统,国家科技创新2030“脑科学与类脑研究”重大项目（2022ZD0208903）
          \item 跨站点多模态脑影像的生成式人工智能方法及诊疗应用研究,国家自然科学基金重点项目（U24A20339）
          \item **** 预警关键技术,国防生物科技优秀青年人才基金项目（****）
          \end{enumerate}

        \fi
    \fi
\else\relax
\fi

\ifisanon\relax
\else
\ifisresumebib

\begingroup
    \settoggle{bbx:gbtype}{false}%局部设置不输出文献类型和载体标识符
	\settoggle{bbx:gbannote}{true}%局部设置输出注释信息
    \renewcommand{\bibfont}{\normalsize}%\fangsong
	%\setcounter{gbnamefmtcase}{6}%局部设置作者的格式为familyahead格式
    \makeatletter
    \renewcommand*{\mkbibnamegiven}[1]{%通过作者注释局部调整作者的格式需与bib配合
    \ifitemannotation{thesisauthor}
    {\ifbibliography{{\textbf{#1}}}{#1}}%
    {#1}\ifbibliography{\ifitemannotation{corresponding}{\textsuperscript{*}}{}}{}%
    }
    \renewcommand*{\mkbibnamefamily}[1]{%
    \ifitemannotation{thesisauthor}
    {\ifbibliography{{\textbf{#1}}}{#1}}
    {#1}}
    \def\blx@maxbibnames{9}
    \def\blx@minbibnames{2}
    %\defcounter{gbbiblocalcase}{1} %局部强迫中文本地化字符串
    %\defcounter{gbbiblocalcase}{2} %局部强迫英文本地化字符串
    %\setlocalbibstring{andotherscn}{et al.}
    %\setlocalbibstring{andothers}{等}
    \renewbibmacro*{author}{%
      \ifboolexpr{
        test \ifuseauthor
        and
        not test {\ifnameundef{author}}
      }
        {\ifisblind%
           \setcounter{numitemval}{1}%
           \whileboolexpr{test {\ifnumcomp{\value{numitemval}}{<}{\value{author}+1}}}
               {\ifitemannotation[author][default][\thenumitemval]{thesisauthor}%
                    {\setcounter{numitemslt}{\thenumitemval}}{}%
               \stepcounter{numitemval}%
               }%
           \ifnumcomp{\value{numitemslt}}{>}{0}
                {\textbf{\iffieldequalstr{userd}{chinese}{第\zhnum{numitemslt}作者}{\Ordinalstring{numitemslt}~Author}}, XXX.}
                {\printnames{author}}
         \else%
            \printnames{author}%
         \fi%%
         \iffieldundef{authortype}
           {}
           {\setunit{\printdelim{authortypedelim}}%
            \usebibmacro{authorstrg}}}
        {}}
    \makeatother


	\begin{refsection}[ref/resume.bib]
	\nocite{zhang2025dmae, zhang2024maser, zhang2022graph, zhang2024mae, wang2024splatflow, liu2024amphibious, su2024m2DC, zhao2025universal, liu2023fusion, lu2022combining, fan2022micro, zhang2023co, cheng2024edgemiformer, zhang2025selection}
    \printbibliography[heading=subbibliography,title={发表的学术论文}]
	\end{refsection}

	\begin{refsection}[ref/resume.bib]
	\nocite{granted-1,  p2024-1, granted-2,p2023-1,  p2024-2, p2025-1, p2023-2, p2025-2}%%
	\printbibliography[heading=subbibliography,title={申请的发明专利}]
	\end{refsection}


	\begin{refsection}[ref/resume.bib]
	\nocite{比赛2024-1, 比赛2024-2, 比赛2023-1, 比赛2023-2, 比赛2023-3, 比赛2023-4, 比赛2022-1}%
	\printbibliography[heading=subbibliography,title={获得的竞赛奖项}]
	\end{refsection}

	\begin{refsection}[ref/resume.bib]
	% \nocite{科技2030, 多模态, 国防2023, 国重2018}%
  \nocite{科技2030, 多模态, 国防2023}%
	\printbibliography[heading=subbibliography,title={参与的科研项目}]
	\end{refsection}
\endgroup


\else

  \section*{发表的学术论文} % 发表的和录用的合在一起

  \begin{enumerate}[label={[\arabic*]},labelsep=6pt,topsep=0pt,partopsep=0pt,
parsep=0pt,itemsep=1pt,leftmargin=0em,itemindent=42pt]
  \item 第一作者.期刊论文.IEEE Transactions on Neural Networks and Learning Systems,2025,SCI影响因子8.9.
  \item 第一作者.期刊论文.IEEE Transactions on Neural Systems and Rehabilitation Engineering,2024,SCI影响因子4.8.
  \item 第一作者.期刊论文.IEEE Transactions on Cognitive and Developmental Systems,2023,SCI影响因子5.0.
  \item 第一作者.会议论文.International Conference on Automation Control,Algorithm,and Intelligent Bionics,2024.
  \item 第二作者.期刊论文.International Journal of Computer Vision,2024,SCI影响因子11.6.
  \item 第二作者.期刊论文.Physics of Fluids,2024,SCI影响因子4.6.
  \item 第四作者.期刊论文.IEEE Transactions on Medical Imaging,2025,SCI影响因子10.6.
  \item 第五作者.期刊论文.Frontiers in Neurorobotics,2025,SCI影响因子2.9.
  \item 第五作者.期刊论文.IEEE Transactions on Human-Machine Systems,2023,SCI影响因子3.5.
  \item 第二作者.会议论文.Youth Academic Annual Conference of Chinese Association of Automation,2022.
  \item 第三作者.会议论文.IEEE International Conference on Civil Aviation Safety and Information Technology,2022.
  \item 第四作者.会议论文.Asian Conference on Artificial Intelligence Technology,2023.
  \item 第二作者.会议论文.Asian Conference on Artificial Intelligence Technology,2025.
  \end{enumerate}

  \section*{申请的发明专利} % 有就写，没有就删除
  \begin{enumerate}[label={[\arabic*]},labelsep=6pt,topsep=0pt,partopsep=0pt,
parsep=0pt,itemsep=1pt,leftmargin=0em,itemindent=42pt]
  \item 第二作者.发明专利.中国,2022,已授权.
  \item 第三作者.发明专利.中国,2022,已授权.
  \item 第二作者.发明专利.中国,2023,实质性审查.
  \item 第二作者.发明专利.中国,2024,实质性审查.
  \item 第二作者.发明专利.中国,2024,实质性审查.
  \item  第二作者.发明专利.中国,2025,实质性审查.
  \item 第三作者.发明专利.中国,2022,实质性审查.
  \item 第二作者.发明专利.中国,2025,已受理.
  \end{enumerate}

  \section*{获得的竞赛奖项} %

\begin{enumerate}[label={[\arabic*]},labelsep=6pt,topsep=0pt,partopsep=0pt,
parsep=0pt,itemsep=1pt,leftmargin=0em,itemindent=42pt]
\normalsize
\item 第一参与人.研究生人工智能创新大赛.国家三等奖,2024.
\item 第一参与人.人工智能技术创新大赛硬件挑战赛.国家二等奖,2023.
\item 第二参与人.人工智能技术创新大赛硬件挑战赛.国家一等奖,2023.
\item 第三参与人.研究生未来飞行器创新大赛.国家二等奖,2023.
\item 第四参与人.研究生机器人创新设计大赛.国家二等奖,2022.
\item 第五参与人.研究生机器人创新设计大赛.国家二等奖,2023.
\item 第五参与人.全国发明展览会技能发展与技术创新大赛.金奖,2024.
\end{enumerate}

\section*{参与的科研项目} %

  \begin{enumerate}[label={[\arabic*]},labelsep=6pt,topsep=0pt,partopsep=0pt,
  parsep=0pt,itemsep=1pt,leftmargin=0em,itemindent=42pt]
  \normalsize
  \item 国家科技创新2030重大项目,2022年-至今.
  \item 国家自然科学基金重点项目,2024年-至今.
  \item 国防生物科技优秀青年人才基金项目,2023年-至今.
  \end{enumerate}

\fi
\fi
\end{resume}


